% GENERATED FILE. Edit canonical lessons, then run npm run generate:books.
% canonical-source-hash: fnv1a64:1709316e9e6fddb7
% canonical-lessons: TE-C151-chhaya, TE-S172-letter-chha, TE-R151-chhaya-recall

\chapter{A Word for Shade — Then Its First Letter}
\label{ch:te-chhaya-chha}
\hlchaptermodality{151}

\begin{chapteropening}
\textbf{By the end of this chapter:} \emph{I can say a Telugu word for shade or shadow, recognise its uncommon first letter, and write that letter in a source-backed order.}

Take \te{ఛ} from the already-known word \te{ఛాయ} and write its five movements in two pen-down runs.
\end{chapteropening}

\section[\texorpdfstring{\te{ఛాయ} (chāya)}{ఛాయ (chāya)}]{\te{ఛాయ} (chāya) — shade, shadow}

\label{lesson:TE-C151-chhaya}



\begin{quote}
Point once to the rare \textbf{\te{ఙ}} you met in \textbf{\te{వాఙ్మయం}}. Now imagine the cool patch beneath a tree. This lesson gives that patch a name, using the familiar \textbf{\te{య}} at the end and a different uncommon letter at the beginning.
\end{quote}

\subsection*{You'll want to know: \te{ఛాయ}}
\textbf{\te{ఛాయ}} — \emph{chāya} — \textquotedblleft{}shade\textquotedblright{} or \textquotedblleft{}shadow\textquotedblright{}.

Say it in two gentle pieces: \emph{chā · ya}. The first sound is an airy \textbf{ch}; the second is the familiar \textbf{\te{య}} sound. Read the word as \textbf{\te{ఛా}-\te{య}}.

\textbf{\te{ఛాయ}} is a learned word with a wider range than one dark shape. In literary or careful language it can also name a reflected image or a tinge of colour. Keep the first meaning small today: shade or shadow.

\begin{grammarlens}[title={What you've built}]
A useful image comes before an unfamiliar letter. You can say the whole word before the next lesson asks your hand to take its first shape apart.
\end{grammarlens}

\subsection*{Your turn}
\begin{itemize}
\raggedright
  \item \emph{Say it:} \emph{chā · ya}, slowly
  \item \emph{Say it:} \emph{chāya}, as one word
  \item \emph{Read it:} \textbf{\te{ఛాయ}} — shade or shadow
  \item \emph{Find:} the familiar \textbf{\te{య}} at the end
\end{itemize}

\begin{tcolorbox}[breakable,colback=teal!4,colframe=teal!35!black,title={Before you move on}]
What does \textbf{\te{ఛాయ}} mean? (\textquotedblleft{}Shade\textquotedblright{} or \textquotedblleft{}shadow.\textquotedblright{}) Say it once more.
\end{tcolorbox}
\section[\texorpdfstring{\te{ఛ} (cha)}{ఛ (cha)}]{\te{ఛ} — the airy first letter in \te{ఛాయ}}

\label{lesson:TE-S172-letter-chha}



\begin{quote}
Say \textbf{\te{ఛాయ}} once. What does it mean? (\textquotedblleft{}Shade\textquotedblright{} or \textquotedblleft{}shadow.\textquotedblright{})

Now slow the first piece: \emph{chā}. The base letter carrying that airy \textbf{ch} is today's whole task.
\end{quote}

\subsection*{Script you'll notice: \te{ఛ}}
\textbf{\te{ఛ}} — \emph{cha}.

In \textbf{\te{ఛాయ}}, the long \textbf{\te{ా}} sign turns the base sound into \textbf{\te{ఛా}}. Remove that vowel sign and the base letter is \textbf{\te{ఛ}}.

This aspirated letter is uncommon in everyday Telugu vocabulary. That is why the book gave you \textbf{\te{ఛాయ}} before asking you to practise the shape.

\subsection*{Writing: \te{ఛ} — five movements, two runs}
\hlblockfigure{\detokenize{figures/TE-S172-letter-chha-filmstrip.pdf}}{How \te{ఛ} is written, stroke by stroke}

Follow the filmstrip slowly. Movements 1–4 stay joined: draw the short upper bar, turn around the left bowl, sweep around the broad outer bowl, and continue through the upper flourish. Lift once. Movement 5 draws the separate stem downward below the middle.

The numbered route is one attested school-style order; Telugu handwriting can vary. Keep the two pen-down runs distinct and let the long curves stay round.

\subsection*{Your turn}
\begin{itemize}
\raggedright
  \item \emph{Trace it:} \textbf{\te{ఛ}} once, following the five numbered movements
  \item \emph{Copy:} \textbf{\te{ఛ}} twice without tracing
  \item \emph{Build:} \textbf{\te{ఛ}} + \textbf{\te{ా}} = \textbf{\te{ఛా}}
  \item \emph{Read it:} \textbf{\te{ఛాయ}}
\end{itemize}

\begin{tcolorbox}[breakable,colback=teal!4,colframe=teal!35!black,title={Before you move on}]
Which letter is this — \textbf{\te{ఛ}}? What word gave it to you? (\textbf{\te{ఛాయ}}.)
\end{tcolorbox}
\section[\texorpdfstring{\te{ఛాయ} (chāya)}{ఛాయ (chāya)}]{\te{ఛాయ} — meaning and shape together}

\label{lesson:TE-R151-chhaya-recall}



\begin{quote}
Picture the cool patch beneath a tree. Say the Telugu word for it. (\textbf{\te{ఛాయ}}.)
\end{quote}

\subsection*{Your turn}
Read \textbf{\te{ఛాయ}} once. Cover it, then write only its first base letter: \textbf{\te{ఛ}}.

Now uncover the word and check two things: the broad curved body and the short separate stem beneath it. Add the long \textbf{\te{ా}} sign only after the base letter is steady.

\begin{itemize}
\raggedright
  \item \emph{Write it:} \textbf{\te{ఛ}} once from memory
  \item \emph{Build:} \textbf{\te{ఛ}} + \textbf{\te{ా}} + \textbf{\te{య}} = \textbf{\te{ఛాయ}}
  \item \emph{Say it:} “shade or shadow”
\end{itemize}

\begin{tcolorbox}[breakable,colback=teal!4,colframe=teal!35!black,title={Before you move on}]
What does \textbf{\te{ఛాయ}} mean, and which base letter begins it? (\textquotedblleft{}Shade\textquotedblright{} or \textquotedblleft{}shadow\textquotedblright{}; \textbf{\te{ఛ}}.)
\end{tcolorbox}
